\documentclass[12pt]{article}
\usepackage{amsmath, amssymb}

\title{CSE 400: Fundamentals of Probability in Computing\\
Lecture Scribe — Tail Bounds and Inequalities}
\date{}

\begin{document}

\maketitle

\section*{Tail and Motivation of Tail Bound}

A tail bound is a bound on the probability that a random variable takes values far from a specified value.

The motivation of tail bounds is to quantify the probability of rare events, particularly events where the random variable deviates significantly from its expectation.

\section*{Markov’s Inequality}

\subsection*{Definition}

Let $X$ be a non-negative random variable. For any $a > 0$:

\[
P(X \ge a) \le \frac{E[X]}{a}
\]

\subsection*{Proof}

Since $X \ge 0$, we write:

\[
E[X] = \sum_x x \cdot P(X = x)
\]

Split the summation into two parts:

\[
E[X] =
\sum_{x < a} x \cdot P(X = x) +
\sum_{x \ge a} x \cdot P(X = x)
\]

For all $x \ge a$, we have $x \ge a$, hence:

\[
\sum_{x \ge a} x \cdot P(X = x) \ge
\sum_{x \ge a} a \cdot P(X = x)
= a \cdot \sum_{x \ge a} P(X = x)
= a \cdot P(X \ge a)
\]

Thus:

\[
E[X] \ge a \cdot P(X \ge a)
\]

Dividing both sides by $a$:

\[
P(X \ge a) \le \frac{E[X]}{a}
\]

\subsection*{Example}

Let $E[X] = 10$. Find an upper bound for $P(X \ge 20)$.

\[
P(X \ge 20) \le \frac{10}{20} = \frac{1}{2}
\]

\section*{Chebyshev’s Inequality}

\subsection*{Definition}

Let $X$ be a random variable with mean $\mu = E[X]$ and variance $\sigma^2 = Var(X)$. For any $k > 0$:

\[
P(|X - \mu| \ge k) \le \frac{\sigma^2}{k^2}
\]

\subsection*{Proof}

Define:

\[
Y = (X - \mu)^2
\]

Since $Y \ge 0$, apply Markov’s inequality:

\[
P(Y \ge k^2) \le \frac{E[Y]}{k^2}
\]

Observe:

\[
P(Y \ge k^2) = P((X - \mu)^2 \ge k^2)
= P(|X - \mu| \ge k)
\]

Also:

\[
E[Y] = E[(X - \mu)^2] = \sigma^2
\]

Thus:

\[
P(|X - \mu| \ge k) \le \frac{\sigma^2}{k^2}
\]

\subsection*{Example}

Let $\mu = 5$, $\sigma^2 = 4$. Find $P(|X - 5| \ge 2)$.

\[
P(|X - 5| \ge 2) \le \frac{4}{2^2} = \frac{4}{4} = 1
\]

\section*{Chernoff Bound and Poisson Tail}

\subsection*{Definition}

For any $t > 0$:

\[
P(X \ge a) = P(e^{tX} \ge e^{ta})
\]

Applying Markov’s inequality:

\[
P(e^{tX} \ge e^{ta}) \le \frac{E[e^{tX}]}{e^{ta}}
\]

Thus:

\[
P(X \ge a) \le \frac{E[e^{tX}]}{e^{ta}}
\]

\subsection*{Proof}

For $t > 0$, the exponential function is monotonic. Therefore:

\[
P(X \ge a) = P(e^{tX} \ge e^{ta})
\]

Applying Markov’s inequality:

\[
P(e^{tX} \ge e^{ta}) \le \frac{E[e^{tX}]}{e^{ta}}
\]

Thus:

\[
P(X \ge a) \le \frac{E[e^{tX}]}{e^{ta}}
\]

\subsection*{Example (Poisson Tail)}

Let $X \sim Poisson(\lambda)$.

\[
E[e^{tX}] = e^{\lambda(e^t - 1)}
\]

Applying Chernoff bound:

\[
P(X \ge a) \le \frac{e^{\lambda(e^t - 1)}}{e^{ta}}
= \exp(\lambda(e^t - 1) - ta)
\]

\section*{Jensen’s Inequality}

\subsection*{Convex Functions}

A function $f$ is convex if for any $x, y$ and $0 \le \theta \le 1$:

\[
f(\theta x + (1 - \theta)y) \le \theta f(x) + (1 - \theta) f(y)
\]

\subsection*{Statement}

Let $X$ be a random variable and $f$ be a convex function. Then:

\[
f(E[X]) \le E[f(X)]
\]

\subsection*{Proof}

Let $X$ take values $x_1, x_2, \ldots, x_n$ with probabilities $p_1, p_2, \ldots, p_n$, where:

\[
\sum_{i=1}^n p_i = 1
\]

Then:

\[
E[X] = \sum_{i=1}^n p_i x_i
\]

By convexity:

\[
f\left(\sum_{i=1}^n p_i x_i\right) \le \sum_{i=1}^n p_i f(x_i)
\]

Thus:

\[
f(E[X]) \le E[f(X)]
\]

\subsection*{Example}

Let $f(x) = x^2$, which is convex.

\[
f(E[X]) = (E[X])^2
\]

\[
E[f(X)] = E[X^2]
\]

Thus:

\[
(E[X])^2 \le E[X^2]
\]

\end{document}